% Usar la plantilla principal: definiciones mínimas y carga del template
\documentclass[spanish,letterpaper,oneside]{article}
% Metadata del documento (ajusta valores según corresponda)
\def\documenttitle {Actividad: Estudio de Caso — Semana 6: Poder, Estado y Derecho}
\def\documentsubtitle {}
\def\documentsubject {Poder, Estado y Derecho}

\def\documentauthor {Nombre del estudiante}
\def\coursename {Filosofía del Derecho}
\def\coursecode {FDE30}

\def\universityname {UNIVERSIDAD ABIERTA Y A DISTANCIA DE MÉXICO}
\def\universityfaculty {Licenciatura en Derecho}
\def\universitydepartment {Programa educativo}
\def\universitylocation {México}
\def\authortable {
	\begin{tabular}{ll}
		Estudiante:
		& \begin{tabular}[t]{l}
			\documentauthor
		\end{tabular} \\
		Materia:
		& \begin{tabular}[t]{l}
			\coursename
		\end{tabular} \\
		Clave:
		& \begin{tabular}[t]{l}
			\coursecode
		\end{tabular} \\
		& \\
		\multicolumn{2}{l}{Fecha de entrega: \today} \\
		\multicolumn{2}{l}{\universitylocation}
	\end{tabular}
}
% Valores adicionales requeridos por la plantilla
\def\universitydepartmentimage {departamentos/UnADM}
\def\universitydepartmentimagecfg {height=1.57cm}

% Cargar la plantilla central usando ruta relativa desde este archivo
\input{template}

\begin{document}

% Generar portada y configuraciones de página según la plantilla
\templatePortrait
\templatePagecfg
\templateIndex
\templateFinalcfg

% ============================================================================
% INTRODUCCIÓN
% ============================================================================

\section*{Introducción}

La relación entre el poder, el Estado y el derecho constituye uno de los núcleos fundamentales del pensamiento jurídico y filosófico contemporáneo. Comprender cómo estas tres dimensiones se interrelacionan es esencial para cualquier profesional del derecho que aspire a ejercer su práctica con ética y fundamentación teórica sólida.

El Estado de derecho no es una realidad automática, sino una construcción conceptual y práctica que requiere la convergencia de principios filosóficos, marcos normativos y mecanismos institucionales. En el caso mexicano, esta estructura se refleja en la Constitución Política de los Estados Unidos Mexicanos, que establece la división de poderes y los límites del ejercicio estatal frente a los derechos de las personas.

El presente análisis examina estos conceptos a través de la bibliografía especializada y el estudio de un caso concreto que permite identificar cómo los principios teóricos se materializan en la práctica jurídica. Este ejercicio es fundamental para desarrollar una visión crítica y reflexiva de los fundamentos del ordenamiento legal mexicano.

\newpage

% ============================================================================
% DEFINICIÓN DE CONCEPTOS
% ============================================================================

\section*{Definición de Conceptos}

\subsection*{Concepto de Estado}

El Estado es la organización política soberana de una comunidad, constituida por una población determinada, un territorio específico y una estructura de poder que ejerce autoridad legítima sobre ese territorio y población. En términos filosóficos, el Estado representa la materialización de la voluntad colectiva y la búsqueda del bien común mediante la aplicación de normas jurídicas. Según \citet{garciaMaynez2002}, el Estado moderno se distingue por su capacidad de monopolizar la coacción legítima y de producir normas vinculantes que regulan la convivencia social.

\subsection*{Elementos del Estado}

Los elementos constitutivos del Estado son tres: (1) \textbf{Población}, que comprende el conjunto de individuos con nacionalidad común que residen en el territorio estatal; (2) \textbf{Territorio}, que abarca el espacio geográfico delimitado donde el Estado ejerce su soberanía, incluyendo tierra, agua y espacio aéreo; y (3) \textbf{Poder soberano}, que es la capacidad del Estado para ejercer autoridad suprema, independiente de cualquier otro poder externo. \citet{ansuategui2007} enfatiza que estos elementos son interdependientes: sin población no hay comunidad política; sin territorio no hay delimitación del ejercicio del poder; sin soberanía no hay capacidad para imponer normas vinculantes.

\subsection*{Formas de Estado}

Las formas de Estado se clasifican según la distribución territorial del poder en: (1) \textbf{Estado unitario}, donde el poder está concentrado en una autoridad central que ejerce jurisdicción en todo el territorio; (2) \textbf{Estado federal}, como México, donde existe una distribución de competencias entre un gobierno central y gobiernos locales (entidades federativas), con autonomía relativa de cada nivel; y (3) \textbf{Estado confederado}, donde los Estados mantienen mayor independencia y delegan ciertos poderes a un órgano común. Mexico es un Estado federal, como se establece en el artículo 40 de la Constitución Política de los Estados Unidos Mexicanos.

\subsection*{Sistema de Gobierno}

El sistema de gobierno se refiere a la estructura de instituciones y órganos mediante los cuales el Estado ejerce su poder y toma decisiones públicas. Los principales sistemas son: (1) \textbf{Sistema presidencial}, caracterizado por la concentración del poder ejecutivo en un presidente electo y la separación de poderes; (2) \textbf{Sistema parlamentario}, donde el parlamento es el órgano central y el ejecutivo depende de la confianza legislativa; y (3) \textbf{Sistema mixto o semi-presidencial}, que combina elementos de ambos. México adopta un sistema presidencial donde el Presidente de la República es tanto jefe de Estado como jefe de gobierno, conforme al artículo 80 constitucional.

\subsection*{División de Poderes}

La división de poderes es el principio fundamental de limitación del poder estatal. Establece que las funciones del Estado deben distribuirse entre tres órganos independientes: (1) \textbf{Poder Legislativo}, encargado de crear normas generales; (2) \textbf{Poder Ejecutivo}, responsable de ejecutar y hacer cumplir las leyes; y (3) \textbf{Poder Judicial}, encargado de interpretar las leyes y resolver controversias. Este principio, desarrollado filosóficamente por Montesquieu, tiene como propósito crear un sistema de \textit{checks and balances} que impida la concentración de poder. En México, se consagra en los artículos 49, 50, 70, 90 y 94 de la Constitución.

\subsection*{Estado de Derecho}

El Estado de derecho es un régimen político donde el ejercicio del poder estatal se encuentra sujeto al ordenamiento jurídico. Sus características fundamentales son: (1) \textbf{Supremacía de la ley}, todas las acciones estatales deben conformarse a la ley; (2) \textbf{Legalidad}, los actos administrativos requieren fundamento legal; (3) \textbf{Protección de derechos fundamentales}, incluyendo derechos humanos; (4) \textbf{Responsabilidad estatal}, los servidores públicos responden por sus actos; y (5) \textbf{Acceso a la justicia}, mediante instituciones independientes. Ansuátegui Roig (2007) sostiene que el Estado de derecho no es únicamente una forma de organización, sino un ideal regulativo que exige la búsqueda continua de la justicia material.

\subsection*{Poder Político, Social y Jurídico}

Aunque frecuentemente se confunden, estos tipos de poder tienen características distintas: (1) \textbf{Poder político} es la capacidad de tomar decisiones vinculantes para toda la comunidad mediante mecanismos institucionales; (2) \textbf{Poder social} es la influencia que ejercen actores no estatales (grupos, asociaciones, empresas) en la dinámica colectiva, sin respaldo coercitivo directo; y (3) \textbf{Poder jurídico} es la facultad de crear, interpretar y aplicar normas mediante órganos especializados. El poder político estatal se distingue por su monopolio legítimo de la coerción, mientras que el poder social opera mediante persuasión e influencia, y el poder jurídico se ejerce en el marco de competencias y procedimientos previamente establecidos.

\newpage

% ============================================================================
% ANÁLISIS DE CASO
% ============================================================================

\section*{Análisis de Caso: Elementos del Estado Mexicano}

El siguiente análisis aplica los conceptos definidos al Estado mexicano, identificando cómo cada elemento fundamental se refleja en la estructura constitucional y en la práctica jurídica contemporánea.

\begin{center}
\begin{tabular}{|m{2.5cm}|m{2.5cm}|m{7.5cm}|}
\hline
\rowcolor[gray]{0.8}
\textbf{Elemento} & \textbf{Artículo} & \textbf{Explicación y Ejemplo} \\
\hline
\textbf{Territorio} & 42 CPEUM & 
El territorio nacional comprende el territorio español según los límites que fija el tratado de Guadalupe Hidalgo de 1848. Entre los límites y fronteras internacionales se encuentra Guatemala y Belice al sur, Estados Unidos al norte. 

\textit{Ejemplo:} Reza García (2024). \textit{Los ríos y cruces fronterizos que dividen a México con Guatemala: Una revisión del tratado de límites.} Demarca la soberanía territorial mediante ríos como el Usumacinta. \\
\hline

\textbf{Población} & 30-34 CPEUM & 
La nacionalidad mexicana es un estatus que se adquiere por nacimiento en territorio mexicano o por descendencia de mexicanos. El artículo 34 define quiénes son ciudadanos mexicanos. La población constituida por ciudadanos mexicanos forma la base política del Estado.

\textit{Ejemplo:} La reforma de 2011 en materia de derechos humanos amplió la definición de quiénes son titulares de derechos y deberes en el Estado mexicano, reconociendo la dignidad de todas las personas. \\
\hline

\textbf{Soberanía} & 39-40 CPEUM & 
La soberanía nacional reside esencialmente en el pueblo (artículo 39). El ejercicio de esta soberanía se realiza mediante los Poderes de la Unión, en los términos establecidos por la Constitución. México es una República representativa, no una democracia directa.

\textit{Ejemplo:} Las elecciones presidenciales cada seis años permiten a la población expresar su voluntad soberana mediante el voto, determinando quién ejercerá el Poder Ejecutivo. \\
\hline

\textbf{División de Poderes} & 49 CPEUM & 
El supremo poder de la Federación se divide para su ejercicio en Legislativo, Ejecutivo y Judicial. Ninguno de estos poderes puede depositarse en un solo órgano, ni puede ejercitarse alguno de ellos por los titulares del otro.

\textit{Ejemplo:} El Congreso (Poder Legislativo) aprueba la Ley de Amparo, el Presidente (Ejecutivo) la promulga, y los Juzgados de Distrito (Judicial) la aplican en casos concretos. Cada uno respeta su competencia. \\
\hline

\textbf{Estado de Derecho} & 16, 17, 1 CPEUM & 
Nadie puede ser molestado en su persona, familia, domicilio, papeles o posesiones, sino por virtud de mandamiento escrito de autoridad competente (artículo 16). La administración pública federal debe actuar conforme a la legalidad, honradez, lealtad, imparcialidad y eficiencia (artículo 1).
asdf
\textit{Ejemplo:} La Ley de Amparo permite a cualquier persona que considere que sus derechos fundamentales han sido violados por una autoridad ejercitar este recurso ante los Juzgados de Distrito, garantizando que aun el Estado está sujeto al derecho. \\
\hline

\textbf{Límites al Poder} & 1 CPEUM & 
En los Estados Unidos Mexicanos todas las personas gozarán de los derechos humanos reconocidos en la Constitución y en los tratados internacionales de los que el Estado es parte. Las normas relativas a derechos humanos se interpretarán de conformidad con la Constitución y los tratados, favoreciendo en todo tiempo a las personas la protección más amplia.

\textit{Ejemplo:} La prohibición de tortura (artículo 22) y el derecho al debido proceso (artículo 20) establecen límites infranqueables al poder punitivo del Estado, incluso en supuestos de criminalidad grave. \\
\hline
\end{tabular}
\end{center}

\newpage

% ============================================================================
% CONCLUSIONES
% ============================================================================

\section*{Conclusiones}

El análisis realizado pone de manifiesto que la relación entre poder, Estado y derecho no es puramente teórica, sino que constituye el fundamento material de la estructura política mexicana. El Estado mexicano, como Estado federal presidencial de derecho, se define precisamente por la subordinación del poder político al ordenamiento jurídico. La Constitución Política de los Estados Unidos Mexicanos establece un marco donde los tres poderes actúan dentro de límites definidos, y donde los derechos fundamentales de las personas representan fronteras que ni siquiera el poder soberano puede traspasar. Esto es lo que distingue el Estado de derecho de una mera dictadura: la existencia de reglas vinculantes para el poder mismo. En el contexto de la práctica profesional del abogado, comprender esta interrelación es fundamental porque permite orientar la acción jurídica hacia la realización de la justicia, no simplemente hacia la ejecución de normas \citep{lovonManualPracticoFilosofia2020}.

Sin embargo, la existencia formal del Estado de derecho en la Constitución no garantiza su realización efectiva. La reforma de 2011 en materia de derechos humanos, la creación de la Comisión Nacional de Derechos Humanos, y los mecanismos como el juicio de amparo, reflejan el reconocimiento institucional de que la brecha entre el deber ser normativo y la realidad fáctica requiere vigilancia constante. La práctica de casos concretos—como los relacionados con igualdad sustantiva entre hombres y mujeres \citep{franzoniacevedoLeyGeneralAcceso2017}, acceso a la justicia \citep{LeyGeneralVictimas}, o prevención de corrupción—demuestra que el Estado de derecho es un proceso continuo, no un logro definitivo. Para el futuro profesional del derecho, esto significa que la labor de un abogado no es únicamente aplicar normas existentes, sino contribuir a la construcción permanente de un Estado donde el poder esté genuinamente subordinado al derecho y a la dignidad humana.

\newpage

% ============================================================================
% REFERENCIAS
% ============================================================================

\section*{Referencias}

\bibliography{referencias/materias/filosofia-derecho/bib-filosofia-derecho/FilosofiaDerecho}

\end{document}
